\documentclass[11pt]{article}
% ===================================================================
% Shared preamble for the "tiny lemma, read slowly" preprint series.
% Mirrors the style of FrobeniusLemmas_preprint.tex.
% ===================================================================

\usepackage{amsmath,amssymb,amsthm}
\usepackage{mathtools}
\usepackage[margin=1.2in]{geometry}
\usepackage{hyperref}
\usepackage{xcolor}
\usepackage{listings}
\usepackage{tikz}
\usepackage{tikz-cd}
\usepackage{mathpartir}
\usepackage{newunicodechar}
\usepackage{setspace}

\onehalfspacing

% ---- Unicode glyphs ----
\newunicodechar{→}{\ensuremath{\to}}
\newunicodechar{↦}{\ensuremath{\mapsto}}
\newunicodechar{↔}{\ensuremath{\leftrightarrow}}
\newunicodechar{⟹}{\ensuremath{\Longrightarrow}}
\newunicodechar{⟶}{\ensuremath{\longrightarrow}}
\newunicodechar{≠}{\ensuremath{\neq}}
\newunicodechar{≤}{\ensuremath{\leq}}
\newunicodechar{≥}{\ensuremath{\geq}}
\newunicodechar{∀}{\ensuremath{\forall}}
\newunicodechar{∃}{\ensuremath{\exists}}
\newunicodechar{∘}{\ensuremath{\circ}}
\newunicodechar{·}{\ensuremath{\cdot}}
\newunicodechar{∈}{\ensuremath{\in}}
\newunicodechar{∉}{\ensuremath{\notin}}
\newunicodechar{≃}{\ensuremath{\simeq}}
\newunicodechar{≅}{\ensuremath{\cong}}
\newunicodechar{≈}{\ensuremath{\approx}}
\newunicodechar{∧}{\ensuremath{\wedge}}
\newunicodechar{∨}{\ensuremath{\vee}}
\newunicodechar{¬}{\ensuremath{\neg}}
\newunicodechar{⊢}{\ensuremath{\vdash}}
\newunicodechar{⟨}{\ensuremath{\langle}}
\newunicodechar{⟩}{\ensuremath{\rangle}}
\newunicodechar{⇑}{\ensuremath{\mathord{\uparrow}}}
\newunicodechar{↑}{\ensuremath{\mathord{\uparrow}}}
\newunicodechar{∎}{\ensuremath{\blacksquare}}
\newunicodechar{∣}{\ensuremath{\mid}}
\newunicodechar{∤}{\ensuremath{\nmid}}
\newunicodechar{∑}{\ensuremath{\textstyle\sum}}
\newunicodechar{∏}{\ensuremath{\textstyle\prod}}
\newunicodechar{∅}{\ensuremath{\emptyset}}
\newunicodechar{∪}{\ensuremath{\cup}}
\newunicodechar{∩}{\ensuremath{\cap}}
\newunicodechar{≡}{\ensuremath{\equiv}}
\newunicodechar{ℕ}{\ensuremath{\mathbb{N}}}
\newunicodechar{ℤ}{\ensuremath{\mathbb{Z}}}
\newunicodechar{ℝ}{\ensuremath{\mathbb{R}}}
\newunicodechar{𝔽}{\ensuremath{\mathbb{F}}}
\newunicodechar{α}{\ensuremath{\alpha}}
\newunicodechar{β}{\ensuremath{\beta}}
\newunicodechar{σ}{\ensuremath{\sigma}}
\newunicodechar{τ}{\ensuremath{\tau}}
\newunicodechar{φ}{\ensuremath{\varphi}}
\newunicodechar{ψ}{\ensuremath{\psi}}
\newunicodechar{χ}{\ensuremath{\chi}}
\newunicodechar{Φ}{\ensuremath{\Phi}}
\newunicodechar{Δ}{\ensuremath{\Delta}}
\newunicodechar{₀}{\textsubscript{0}}
\newunicodechar{₁}{\textsubscript{1}}
\newunicodechar{₂}{\textsubscript{2}}
\newunicodechar{ᵏ}{\ensuremath{{}^{k}}}
\newunicodechar{ⁿ}{\ensuremath{{}^{n}}}
\newunicodechar{ʲ}{\ensuremath{{}^{j}}}
\newunicodechar{²}{\ensuremath{{}^{2}}}
\newunicodechar{³}{\ensuremath{{}^{3}}}
\newunicodechar{⁴}{\ensuremath{{}^{4}}}
\newunicodechar{⁻}{\ensuremath{{}^{-}}}
\newunicodechar{¹}{\ensuremath{{}^{1}}}

% ---- Lean listing style ----
\definecolor{leanbg}{RGB}{247,247,250}
\definecolor{leankw}{RGB}{0,90,160}
\definecolor{leancm}{RGB}{120,120,120}
\lstdefinestyle{lean}{
  backgroundcolor=\color{leanbg},
  basicstyle=\ttfamily\small,
  keywordstyle=\color{leankw}\bfseries,
  commentstyle=\color{leancm}\itshape,
  morekeywords={theorem,lemma,def,fun,Prop,Type,variable,where,
    Field,Fintype,Finite,DecidableEq,CommRing,CommSemiring,Ring,CharP,Function,
    Injective,Surjective,Bijective,by,rfl,have,rw,exact,ext,simp,intro,
    iterateFrobenius,RingHom,Commute,convert,apply,using,Nat,Coprime,gcd,
    Odd,Even,IsAPN,IsAB,walsh,Nat,obtain,induction,cases,calc,grind},
  showstringspaces=false,
  columns=fullflexible,
  extendedchars=true,
  frame=single,
  framesep=6pt,
  xleftmargin=4pt,
  aboveskip=12pt,
  belowskip=14pt,
}
\lstset{style=lean}

\newtheorem{lemma}{Lemma}
\newtheorem{theorem}{Theorem}
\newtheorem{corollary}{Corollary}

% boxed "what just happened" note
\newcommand{\note}[1]{%
  \par\smallskip
  \noindent\colorbox{leanbg}{\parbox{\dimexpr\linewidth-2\fboxsep}{\small #1}}%
  \par\smallskip}


\title{\textbf{The glue exponent}\\[4pt]
\large $d(k)\cdot(2^k+1)=2^{3k}+1$, read slowly}
\author{}
\date{}

\begin{document}
\maketitle

\begin{abstract}
\noindent
One arithmetic identity.\\[2pt]
$\bullet$\quad \texttt{d\_mul\_gold}:\;\;
$d(k)\cdot(2^k+1)=2^{3k}+1$,\quad where $d(k)=2^{2k}-2^k+1$.\\[6pt]
\noindent
Two lines of \texttt{zify}/\texttt{ring}, but it is the algebraic glue between the
\emph{Kasami} exponent $d(k)$ and the \emph{Gold} exponent $2^k+1$, and it underlies
``decompose $\Phi$ as a composition of two bijections''. The cross-form
factorization plays the same connective role.
\end{abstract}

\bigskip

% =====================================================================
\section*{Two exponents}

\begin{itemize}\itemsep4pt
  \item \textbf{Gold:}\;\; $2^k+1$.
  \item \textbf{Kasami:}\;\; $d(k)=2^{2k}-2^k+1$.
\end{itemize}
They look unrelated. The identity says they multiply to a clean cube:
\[
d(k)\cdot(2^k+1)\;=\;2^{3k}+1 .
\]

\bigskip

% =====================================================================
\section{The lemma}

\begin{lstlisting}
theorem d_mul_gold (k : ℕ) (hk : k ≥ 1) :
    d k * (2 ^ k + 1) = 2 ^ (3 * k) + 1 := by
  unfold d
  zify
  rw [Nat.cast_sub (by gcongr <;> linarith)]
  push_cast
  ring
\end{lstlisting}

\note{\texttt{d} is defined with \emph{truncated} natural subtraction
($2^{2k}-2^k$), so the proof first lifts to $\mathbb{Z}$ with \texttt{zify} and
\texttt{Nat.cast\_sub}, where subtraction behaves, then finishes with \texttt{ring}.}

\bigskip
\subsection*{Why it is true (the cube identity)}

\noindent
Over $\mathbb{Z}$, set $q=2^k$. Then $d(k)=q^2-q+1$ and the Gold factor is $q+1$.
This is the classical sum-of-cubes factoring:
\[
(q+1)(q^2-q+1) \;=\; q^3+1 \;=\; 2^{3k}+1 .
\]

\[
\begin{tikzcd}[column sep=large]
\underbrace{2^k+1}_{\text{Gold}} \;\times\; \underbrace{2^{2k}-2^k+1}_{\text{Kasami } d(k)}
\arrow[r, equal, "{q^3+1}"]
& 2^{3k}+1
\end{tikzcd}
\]

\note{$x^3+1=(x+1)(x^2-x+1)$ is the reason the Gold exponent and the Kasami
exponent are ``complementary cofactors'' of the cube $2^{3k}+1$.}

\bigskip

% =====================================================================
\section{Where it lands: composing two bijections}

\noindent
On $\mathrm{GF}(2^n)$, raising to a power is a bijection iff the exponent is
coprime to $2^n-1$. The identity factors the Kasami power map through Gold:
\[
\bigl(x^{\,d(k)}\bigr)^{\,2^k+1} \;=\; x^{\,d(k)\,(2^k+1)} \;=\; x^{\,2^{3k}+1}.
\]
So the cube power $x\mapsto x^{2^{3k}+1}$ \emph{is} Gold-after-Kasami. When both
right-hand factors are permutations, the identity lets the analysis
\emph{decompose} the differential as a composition of two bijections --- the
strategy behind the Dempwolff--M\"uller route to Kasami APN.

\bigskip
\subsection*{The same role: cross-form factorization}

\noindent
The companion connective identity is the cross form. With
$L_k(P)=P^{2^k}+P$ and $N_k(s)=s^{2^k+1}$:
\[
\mathrm{Cross}_k(s,P)\;=\;s\,P^{2^k}+s^{2^k}P
\;=\; N_k(s)\cdot L_k\!\bigl(P/s\bigr).
\]

\begin{lstlisting}
def L (k : ℕ) (x : F) : F := x ^ (2 ^ k) + x
def N (k : ℕ) (x : F) : F := x ^ (2 ^ k + 1)
def Cross (k : ℕ) (s P : F) : F := s * P ^ (2 ^ k) + s ^ (2 ^ k) * P
\end{lstlisting}

\noindent
Just as \texttt{d\_mul\_gold} factors an \emph{exponent} as Gold $\times$ Kasami,
the cross form factors a \emph{value} as Gold-norm $N_k(s)$ times truncated-trace
$L_k(P/s)$. Both turn one hard object into a product of two understood ones.

\note{Pattern: a one-line algebraic factoring (of an exponent, or of a value)
turns ``analyse $\Phi$'' into ``analyse a product of two bijections''. The cube
identity $x^3+1=(x+1)(x^2-x+1)$ and the cross factorization are the two faces of
this glue.}

\bigskip

% =====================================================================
\section*{One-screen summary}

\begin{center}
\renewcommand{\arraystretch}{1.6}
\begin{tabular}{l l}
\textbf{object} & \textbf{factoring} \\
\hline
exponent & $d(k)\,(2^k+1)=2^{3k}+1$ \;($q^3+1=(q+1)(q^2-q+1)$) \\
power map & $x^{2^{3k}+1}=(x^{d(k)})^{2^k+1}$ \;(Gold $\circ$ Kasami) \\
value & $\mathrm{Cross}_k(s,P)=N_k(s)\cdot L_k(P/s)$ \\
\end{tabular}
\end{center}

\bigskip
\noindent
\colorbox{leanbg}{\parbox{\dimexpr\linewidth-2\fboxsep}{\centering
$d(k)\,(2^k+1)=2^{3k}+1$ $\;\Rightarrow\;$ Kasami power $=$ Gold $\circ$ Kasami
$\;\Rightarrow\;$ ``compose two bijections''.}}

\end{document}
